\pdfoutput=1
\documentclass[11pt]{article}
\usepackage[T1]{fontenc}
\usepackage{lmodern}
\usepackage{amsmath,amssymb,amsthm,mathtools}
\usepackage{booktabs}
% Generated by scripts/make_claims_tex.py from claims.json -- do not edit by hand.
\newcommand{\clSigmaElo}{3.8724615957722}
\newcommand{\clSigmaEhi}{3.8724615957734}
\newcommand{\clareaElo}{0.8494076787253}
\newcommand{\clareaEhi}{0.8494076787257}
\newcommand{\cletaMax}{1.9559\times10^{-13}}
\newcommand{\clsigmaTilde}{1.290820531924}
\newcommand{\clsigmaThreeLo}{2.2003}
\newcommand{\clcrNh}{16}
\newcommand{\clcrDofs}{408}
\newcommand{\clcrMu}{3}
\newcommand{\clcrNeg}{3}
\newcommand{\clcrChSq}{0.121137}
\newcommand{\clrhoOne}{0.6}
\newcommand{\clpsiMax}{6.3}
\newcommand{\cllocalBoxes}{460}
\newcommand{\clepsBdryMax}{2.71\times10^{-13}}
\newcommand{\clepsGradMax}{4.01\times10^{-13}}
\newcommand{\clconeEig}{1.857812}
\newcommand{\clgramE}{1.577442}
\newcommand{\clconeSlope}{3.53321}
\newcommand{\clfarIntervals}{16}
\newcommand{\clfarMax}{3.1144}
\newcommand{\clfarBound}{3.12}
\newcommand{\clcrSigmaEight}{1.28042}
\newcommand{\clcrSigmaSixteen}{1.28790}
\newcommand{\clthetaZero}{1.83525}
\newcommand{\clthetaA}{1.93168}
\newcommand{\clthetaB}{2.40957}
\newcommand{\clthetaC}{2.89159}
\newcommand{\clthinAgree}{4\times10^{-4}}
\newcommand{\clTequiEst}{3.8724615957728}
\newcommand{\clTequiErr}{1\times10^{-13}}
\newcommand{\clTequiUp}{3.8724615957729}
\newcommand{\clTfiftyEst}{2.87773098574}
\newcommand{\clTfiftyErr}{3\times10^{-11}}
\newcommand{\clTfiftyUp}{2.8777309857377}
\newcommand{\clTseventyEst}{2.573076572892}
\newcommand{\clTseventyErr}{2\times10^{-12}}
\newcommand{\clTseventyUp}{2.5730765728918}
\newcommand{\clTfortyEst}{2.53503581509}
\newcommand{\clTfortyErr}{2\times10^{-11}}
\newcommand{\clTfortyUp}{2.5350358150935}
\newcommand{\clTrightisoEst}{2.4960723700475}
\newcommand{\clTrightisoErr}{1\times10^{-13}}
\newcommand{\clTrightisoUp}{2.4960723700476}
\newcommand{\clThalfeqEst}{1.9013817577255}
\newcommand{\clThalfeqErr}{1\times10^{-13}}
\newcommand{\clThalfeqUp}{1.9013817577256}
\newcommand{\clTobtuseEst}{1.5603441}
\newcommand{\clTobtuseErr}{3\times10^{-7}}
\newcommand{\clTobtuseUp}{1.5603440800120}
\newcommand{\clTsquareEst}{5.5060219386901}
\newcommand{\clTsquareErr}{1\times10^{-13}}
\newcommand{\clTsquareUp}{5.5060219386902}
\newcommand{\clTpentagonEst}{5.9774726}
\newcommand{\clTpentagonErr}{4\times10^{-7}}
\newcommand{\clTpentagonUp}{5.9774725550439}
\newcommand{\clThexagonEst}{6.135570}
\newcommand{\clThexagonErr}{5\times10^{-6}}
\newcommand{\clThexagonUp}{6.1355701455936}

\usepackage[margin=1in]{geometry}
\usepackage[colorlinks=true,urlcolor=blue,citecolor=blue]{hyperref}

\newtheorem{theorem}{Theorem}[section]
\newtheorem{proposition}[theorem]{Proposition}
\newtheorem{lemma}[theorem]{Lemma}
\newtheorem{corollary}[theorem]{Corollary}
\theoremstyle{definition}
\newtheorem{definition}[theorem]{Definition}
\theoremstyle{remark}
\newtheorem{remark}[theorem]{Remark}
\newtheorem{observation}[theorem]{Observation}

\emergencystretch=3em
\newcommand{\R}{\mathbb{R}}
\newcommand{\C}{\mathbb{C}}
\newcommand{\eps}{\varepsilon}
\newcommand{\dd}{\,\mathrm{d}}
\newcommand{\Lam}{\Lambda}
\DeclareMathOperator{\tr}{tr}
\renewcommand{\Re}{\operatorname{Re}}
\renewcommand{\Im}{\operatorname{Im}}
\DeclareMathOperator{\sinc}{sinc}

\title{The equilateral triangle maximises the first Steklov eigenvalue among triangles}
\author{Jason Davies\thanks{\href{mailto:jason@jasondavies.com}{\nolinkurl{jason@jasondavies.com}}}}
\date{November 19, 2026}

\begin{document}

\maketitle
\begin{abstract}
Let $\sigma_1(T)$ be the first nonzero Steklov eigenvalue of a triangle $T$ and $L(T)$ its perimeter.
Among triangles the equilateral triangle $E$ uniquely maximises $\sigma_1 L$. If the shortest and longest
sides of $T$ have lengths $a\le b$, then
$\sigma_1(E)L(E)-\sigma_1(T)L(T)\ge (b-a)/(2b)$.
Since $|T|/L(T)^2$ is also maximal at $E$, the equilateral triangle uniquely maximises
$\sigma_1\sqrt{|T|}$ as well. This settles the case $n=3$ of the Steklov analogue of the P\'olya--Szeg\H{o}
conjecture for $n$-gons, for both normalisations. The extremal value satisfies
$3.8724615957722\le \sigma_1(E)L(E)\le 3.8724615957734$. For side length one,
$\sigma_1(E)=\sigma_2(E)<\sigma_3(E)=2\sqrt3<\sigma_4(E)$.
Under displacement of the apex over a fixed unit base, the first two perimeter-normalised eigenvalues
split linearly with equal and opposite slopes independent of direction.

The proof is computer-assisted. Near $E$ we transplant the first eigenspace of $E$ affinely; away from $E$
the trial function $\cos(\pi x)$ suffices. We verify the finite matrix and interval inequalities in ball arithmetic.
\end{abstract}

\noindent\textbf{AI transparency statement.} The main result, its
computer-assisted proof, the partial Lean formalisation and the manuscript
were produced by artificial intelligence. The human author initiated the
broader research and directed the AI systems, but did not write the text
and has not checked the arguments.
Details are given at the end of the paper.

\section{Introduction}

For a bounded Lipschitz domain $\Omega\subset\R^2$, the Steklov problem
$\Delta u=0$ in $\Omega$, $\partial_\nu u=\sigma u$ on $\partial\Omega$
has eigenvalues $0=\sigma_0<\sigma_1(\Omega)\le\sigma_2(\Omega)\le\cdots$. They satisfy
\begin{equation}\label{eq:var}
\sigma_1(\Omega)=\min_{u\in H^1(\Omega)\setminus\R}\mathcal R_\Omega(u),\qquad
\mathcal R_\Omega(u)=\frac{\int_\Omega|\nabla u|^2}{\int_{\partial\Omega}u^2-\frac1{|\partial\Omega|}\bigl(\int_{\partial\Omega}u\bigr)^2}.
\end{equation}
The scale-invariant quantity is
\[\Sigma(\Omega):=\sigma_1(\Omega)\,L(\Omega),\qquad L(\Omega)=|\partial\Omega|.\]
Weinstock's inequality~\cite{Weinstock} states that $\Sigma(\Omega)\le2\pi$ for simply connected planar
domains, with equality only for the disk; see the survey~\cite{GP}. For polygons with a prescribed number of
sides, the natural conjecture is the analogue of P\'olya--Szeg\H{o}'s conjecture for the Dirichlet Laplacian:

\medskip\noindent\textbf{Problem} (Oudet; AIM problem list \emph{Steklov eigenproblems}, Problem 4.6~\cite{AimPL}).
\emph{Do regular $n$-gons maximise $\sigma_1$ among $n$-gons of given area, or of given perimeter?}

\medskip
Dominguez, Nigam and Shahriari~\cite{DNS} gave numerical evidence for triangles using finite elements and
Bayesian optimisation. Nigam also posed the perimeter question as Open Problem~8 at the
2020 miniconference on sharp eigenvalue estimates~\cite{Nigam}. Among rectangles of given perimeter or area, the square maximises $\sigma_1$:
this is the planar case of the cuboid inequality of Girouard, Lagac\'e, Polterovich and
Savo~\cite[Thm.~1.6]{GLPS}. Cheng, Gui, Hu, Li and Yao~\cite{CGHLY} state the conjecture for convex polygons
with at most $N$ sides at fixed perimeter (their Conjecture~1.1), and leave it open for every $N\ge3$. They
prove that $\sigma_1$ of the perimeter-normalised regular $N$-gon increases strictly with $N$, and that its
first eigenspace is two-dimensional. Dryden, Gordon, Moreno, Rowlett and Villegas-Blas~\cite{DGMRV} prove
$\sigma_k(T)L(T)\le\frac{8\sqrt3}{3}\pi^2k^2\vartheta(T)$ for triangles, where $\vartheta(T)$ is the smallest
angle. This shows that degenerate triangles are far from optimal, but it does not compare triangles with $E$.

For the Neumann problem, Laugesen and Siudeja~\cite{LS} proved that the equilateral triangle maximises
the first nonzero Neumann eigenvalue times $L^2$ among triangles, and hence times the area. Near the
equilateral triangle their trial functions are the eigenfunctions of the equilateral triangle, transplanted
by the linear map onto the given triangle. This affine transplantation is the precedent for our local
argument. The Steklov problem adds two difficulties: the first eigenvalue of $E$ is double, so the maximum
is a cone and the comparison degenerates at $E$; and a rigorous lower bound for $\sigma_1(E)$ is needed. We are
not aware of a proof of the Steklov case for any $n\ge3$.

\begin{theorem}\label{thm:main}
Let $E$ be an equilateral triangle. For every triangle $T$,
\[\Sigma(T)\le\Sigma(E),\]
with equality if and only if $T$ is equilateral. More precisely, normalise a longest side to $[0,1]$ and
write the apex as $(p,q)$ with $p\le1/2$ and $q>0$. If $\rho=|(p,q)-(1/2,\sqrt3/2)|$, then
\[
\Sigma(E)-\Sigma(T)\ \ge\ \frac{\rho}{2}\ \ge\
\frac{\ell_{\max}-\ell_{\min}}{2\ell_{\max}},
\]
where $\ell_{\max}$ and $\ell_{\min}$ are the extreme side lengths of $T$. Moreover
\[\clSigmaElo\le\Sigma(E)\le\clSigmaEhi .\]
\end{theorem}

Girouard and Polterovich computed the first sixty normalised eigenvalues of the equilateral triangle
numerically~\cite[Example~3.2.1]{GP}. For side length one we obtain $\sigma_1(E)=\sigma_2(E)<\sigma_3(E)=2\sqrt3<\sigma_4(E)$
(Proposition~\ref{prop:eq}). The double eigenvalue agrees with~\cite[Thm.~1.2]{CGHLY} for $N=3$.

\begin{corollary}\label{cor:area}
For every triangle $T$,
\[\sigma_1(T)\sqrt{|T|}\le\sigma_1(E)\sqrt{|E|}\in[\clareaElo,\,\clareaEhi],\]
with equality if and only if $T$ is equilateral.
\end{corollary}

\begin{proof}
$\sigma_1\sqrt{|T|}=\Sigma(T)\,(|T|/L^2)^{1/2}$. By Heron's formula and the AM--GM inequality,
$|T|/L^2\le\sqrt3/36$, with equality only for equilateral triangles. Now apply Theorem~\ref{thm:main}.
\end{proof}

\subsection*{Method}
Upper bounds for $\sigma_1(T)$ come directly from~\eqref{eq:var}; a lower bound is needed only for the
single domain $E$. We parametrise triangles by the apex $(p,q)$ over the base $[0,1]$
(Lemma~\ref{lem:param}), and let $\rho$ be the distance of $(p,q)$ from the equilateral apex.
\begin{enumerate}
\item \emph{Near $E$} ($\rho\le\rho_1\approx0.6$). The trial space is the exact first eigenspace of $E$, transplanted by
the affine map $E\to T$. The comparison involves only a few moments of the eigenspace and two exact
identities. These identities let us divide out the factor $\rho$ (Lemma~\ref{lem:XrhoY}), which turns the
degenerate comparison at $\rho=0$ into a sign condition. We check the sign condition on boxes in ball
arithmetic. The moments are enclosed using a harmonic-polynomial approximation, a residual bound for the
Dirichlet-to-Neumann map, and spectral isolation of a cubic eigenfunction by a Crouzeix--Raviart bound for $\sigma_4(E)$.
\item \emph{Away from $E$} ($\rho\ge\rho_1$). The single trial function $\cos(\pi x)$ gives
$\Sigma(T)\le\clfarBound$, well below $\Sigma(E)>3.8724$. A one-dimensional check on $\clfarIntervals$ intervals suffices.
\end{enumerate}

\section{Preliminaries}

\subsection{The Dirichlet-to-Neumann map}
Let $\Omega$ be a bounded Lipschitz domain. The Dirichlet-to-Neumann map $\Lam$ is the self-adjoint
operator on $L^2(\partial\Omega)$ associated with the closed form $f\mapsto\int_\Omega|\nabla\mathcal Hf|^2$ on
$H^{1/2}(\partial\Omega)$, where $\mathcal Hf$ is the harmonic extension. It has compact resolvent and eigenvalues
$\sigma_k$, with an orthonormal basis of eigenvectors $\chi_k$ of $L^2(\partial\Omega)$; $\chi_0$ is
constant. If $u$ is a harmonic polynomial with trace $f$, Green's formula gives $f\in\operatorname{dom}\Lam$ and
$\Lam f=\partial_\nu u$. For $\tilde\sigma\in\R$ and $\eta>0$ write $J=[\tilde\sigma-\eta,\tilde\sigma+\eta]$.

\begin{lemma}[Inclusion]\label{lem:incl}
Let $W\subset\operatorname{dom}\Lam$ have dimension $k$, and suppose $\|(\Lam-\tilde\sigma)f\|\le\eta\|f\|$ for all
$f\in W$. Then $\Lam$ has at least $k$ eigenvalues, counted with multiplicity, in $J$.
\end{lemma}
\begin{proof}
Suppose not. Then some $f\in W\setminus0$ is orthogonal to the spectral subspace of $J$, and
$\|(\Lam-\tilde\sigma)f\|^2=\sum_{|\sigma_j-\tilde\sigma|>\eta}|\sigma_j-\tilde\sigma|^2|\langle f,\chi_j\rangle|^2>\eta^2\|f\|^2$,
a contradiction.
\end{proof}

\begin{lemma}[Eigenvector enclosure; cf.~\cite{DK}]\label{lem:enclosure}
Let $P$ be the spectral projection of $\Lam$ onto $J$. Suppose every eigenvalue of $\Lam$ outside $J$ is
either $0$ or at least $\sigma^{\rm lo}>\tilde\sigma\ge0$, and put $\gamma=\sigma^{\rm lo}-\tilde\sigma$. Let
$f\in\operatorname{dom}\Lam$ be orthogonal to the constants, and $r=(\Lam-\tilde\sigma)f$. Then
\[\|(I-P)f\|_{L^2(\partial\Omega)}\le\frac{\|r\|}{\gamma},\qquad
\int_\Omega\bigl|\nabla\mathcal H(I-P)f\bigr|^2\le\frac{\sigma^{\rm lo}}{\gamma^2}\|r\|^2.\]
\end{lemma}
\begin{proof}
$(\sigma_j-\tilde\sigma)\langle f,\chi_j\rangle=\langle r,\chi_j\rangle$ for every $j$, and the component along
$\chi_0$ vanishes. The remaining $j$ with $\sigma_j\notin J$ have $\sigma_j\ge\sigma^{\rm lo}$, so
$\|(I-P)f\|^2\le\|r\|^2/\gamma^2$, and
$\int|\nabla\mathcal H(I-P)f|^2=\sum\sigma_j|\langle r,\chi_j\rangle|^2/(\sigma_j-\tilde\sigma)^2\le\sigma^{\rm lo}\|r\|^2/\gamma^2$,
because $\sigma/(\sigma-\tilde\sigma)^2$ is decreasing for $\sigma>\tilde\sigma$ when $\tilde\sigma\ge0$ (its derivative is $-(\sigma+\tilde\sigma)/(\sigma-\tilde\sigma)^3$).
\end{proof}

\subsection{Parametrisation}
For $(p,q)$ with $q>0$, write
\[T(p,q)=\operatorname{conv}\{(0,0),(1,0),(p,q)\},\]
and set
\[F=\{(p,q):0<p\le\tfrac12,\ q>0,\ (1-p)^2+q^2\le1\}.\]
Since $\Sigma$ is invariant under similarities, it suffices to treat the triangles $T(p,q)$, $(p,q)\in F$.

\begin{lemma}\label{lem:param}
Every triangle is similar to $T(p,q)$ for some $(p,q)\in F$, and $T(p,q)$ is equilateral if and only if
$(p,q)=(\frac12,\frac{\sqrt3}{2})$.
\end{lemma}
\begin{proof}
Rescale so that a longest side has length $1$. Move it to $[0,1]\times\{0\}$ with the third vertex
$(p,q)$ in the upper half-plane, and reflect in $x=\frac12$ if necessary so that $p\le\frac12$. The side
from $(1,0)$ to the apex is not longer than the base, so $(1-p)^2+q^2\le1$; this forces $p>0$. The other
side then satisfies $p^2+q^2\le(1-p)^2+q^2\le1$. The equilateral case is clear.
\end{proof}

We write $p=\frac12+\rho\cos\psi$, $q=\frac{\sqrt3}2+\rho\sin\psi$ for polar coordinates about the
equilateral apex; $\rho$ is the Euclidean distance in the $(p,q)$-plane.

\subsection{Affine transplantation}
Let $A\in GL_2^+(\R)$, $u\in H^1(\Omega)$ and $u_A=u\circ A^{-1}\in H^1(A\Omega)$. Write $G=A^{-1}A^{-T}$. For
an edge $e$ of a polygon $\Omega$ with unit tangent $\tau_e$, put $s_e=|A\tau_e|$. Changing variables gives
\begin{equation}\label{eq:transplant}
\begin{gathered}
\int_{A\Omega}|\nabla u_A|^2=\det A\sum_{i,j}G_{ij}\int_\Omega\partial_iu\,\partial_ju,\\
\int_{Ae}u_A^2=s_e\int_eu^2,\qquad \int_{Ae}u_A=s_e\int_eu,
\end{gathered}
\end{equation}
because the arc length on $Ae$ is $s_e$ times that on $e$. For a finite-dimensional space $\mathcal V$ of trial
functions, define the forms
$\mathcal N_A(u,u')=\int\nabla u_A\cdot\nabla u'_A$ and
$\mathcal D_A(u,u')=\sum_es_e\int_euu'-\frac1{L(A\Omega)}(\sum_es_e\int_eu)(\sum_es_e\int_eu')$ on $\mathcal V$.
If $\mathcal D_A$ is positive definite on $\mathcal V$, then~\eqref{eq:var} gives
$\sigma_1(A\Omega)\le\min_{\mathcal V\setminus 0}\mathcal N_A/\mathcal D_A$.

\section{The equilateral triangle}\label{sec:eq}

Let $E=T(\frac12,\frac{\sqrt3}2)$, with side $1$ and $L(E)=3$. Let $z_0$ be its centroid and
$w=-i\sqrt3\,(z-z_0)$, so that the circumradius is $1$ in $w$ and the apex lies on the positive real $w$-axis. The dihedral group $D_3$
acts by $w\mapsto e^{2\pi i/3}w$ and $w\mapsto\bar w$.

\subsection{An exactly equivariant trial pair}
Given real coefficients $c_k$ ($k\le20$, $k\not\equiv0\bmod3$, fourteen in all), put
$\mathfrak f=\sum_{k\equiv1}c_kw^k+\sum_{k\equiv2}c_k\bar w^k$, $u_1=\Re\mathfrak f$ and $u_2=\Im\mathfrak f$. Rotation by
$2\pi/3$ multiplies $\mathfrak f$ by $e^{2\pi i/3}$, and reflection conjugates it. Hence $(u_1,u_2)$ transforms
under $D_3$ exactly as the coordinates $(\Re w,\Im w)$ do, whatever the $c_k$ are. In particular
$\int_{\partial E}u_a=0$; each $u_a$ is the real part of a holomorphic polynomial $\mathfrak f_a$ (since
$\Re\bar w^k=\Re w^k$ and $\Im\bar w^k=-\Im w^k$), and by Schur's lemma all their Gram
matrices are scalar. The $c_k$ come from a floating-point Ritz computation and are then fixed as exact
binary numbers.

Let $\tilde\sigma\approx\clsigmaTilde$ be the Rayleigh quotient of $u_1$, rounded to binary64. All
boundary integrals of products of $u_a$ and $\partial_\nu u_a$ are integrals of polynomials along segments,
and we compute them exactly in ball arithmetic. This gives the relative residual
\[\eta:=\max_{f\in\operatorname{span}(u_1,u_2)\setminus0}\frac{\|\partial_\nu f-\tilde\sigma f\|_{L^2(\partial E)}}{\|f\|_{L^2(\partial E)}}\le\cletaMax.\]
By Lemma~\ref{lem:incl}, $\Lam_E$ has at least two eigenvalues in $J=[\tilde\sigma-\eta,\tilde\sigma+\eta]$.

\subsection{\texorpdfstring{A Crouzeix--Raviart lower bound for $\sigma_4(E)$}{A Crouzeix-Raviart lower bound for sigma4(E)}}
Subdivide $E$ uniformly into $n_h^2$ equilateral cells $K$ of side $h=1/n_h$. Let $\mathcal C_h$ be the
Crouzeix--Raviart space~\cite{CR}: piecewise linear functions that are continuous at the midpoints $x_e$ of
interior edges. Let $\Pi_h$ be the interpolation by edge means, $(\Pi_hu)(x_e)=|e|^{-1}\int_eu$. Write
$\mathfrak a_h(u,y)=\sum_K\int_K\nabla u\cdot\nabla y$ and $\mathfrak b(u,y)=\int_{\partial E}uy$, and let
$\sigma_{h,0}\le\sigma_{h,1}\le\cdots$ be the finite eigenvalues of the pencil $(\mathfrak a_h,\mathfrak b)$
on $\mathcal C_h$. The min--max principle holds over subspaces on which $\mathfrak b$ is positive definite.

\begin{lemma}[Trace constant]\label{lem:trace}
Let $K$ be an equilateral triangle of side $h$, $e$ an edge of $K$, and $\varphi\in H^1(K)$ with
$\int_e\varphi=0$. Then
\[\|\varphi\|^2_{L^2(e)}\le\frac{4}{\sqrt3}\Bigl(\frac1{\pi^2}+\frac1\pi\Bigr)h\,\|\nabla\varphi\|^2_{L^2(K)}.\]
\end{lemma}
\begin{proof}
Let $\hat\varphi=\varphi-\bar\varphi_K$. Since $\varphi$ has mean zero on $e$, $\|\varphi\|_e\le\|\varphi-\zeta\|_e$ for every
constant $\zeta$; in particular $\|\varphi\|_e\le\|\hat\varphi\|_e$. Let $O_e$ be the vertex opposite $e$. The
field $(x-O_e)\hat\varphi^2$ has normal component $0$ on the two edges through $O_e$, and $2|K|/|e|$ times
$\hat\varphi^2$ on $e$. Its divergence is $2\hat\varphi^2+2\hat\varphi\nabla\hat\varphi\cdot(x-O_e)$. The divergence
theorem therefore gives
\[\|\hat\varphi\|_e^2=\frac{|e|}{|K|}\Bigl(\|\hat\varphi\|_K^2+\int_K\hat\varphi\,\nabla\hat\varphi\cdot(x-O_e)\Bigr)
\le\frac{|e|}{|K|}\bigl(\|\hat\varphi\|_K^2+h\|\hat\varphi\|_K\|\nabla\varphi\|_K\bigr).\]
By the Payne--Weinberger inequality~\cite{PW,Beb} on the convex set $K$ of diameter $h$,
$\|\hat\varphi\|_K\le\frac h\pi\|\nabla\varphi\|_K$. Finally $|e|/|K|=4/(\sqrt3h)$.
\end{proof}

\begin{lemma}[Projection bound; cf.~\cite{Liu,YXL}]\label{lem:liu}
Let $n_h\ge2$ and $C_h^2=2\cdot\frac{4}{\sqrt3}(\frac1{\pi^2}+\frac1\pi)h$. Then for each $k$ below the number of
boundary edges,
\[\sigma_k(E)\ge\frac{\sigma_{h,k}}{1+C_h^2\sigma_{h,k}}.\]
\end{lemma}
\begin{proof}
(i) For every $y_h\in\mathcal C_h$, $\mathfrak a_h(u-\Pi_hu,y_h)=0$. Indeed $\nabla y_h$ is constant on each
cell $K$, and $\int_K\nabla(u-\Pi_hu)=\int_{\partial K}(u-\Pi_hu)\nu=0$ because $\Pi_hu|_K$ is linear and has
the same edge means as $u$.

(ii) Write $\varphi=u-\Pi_hu$. It has mean zero on every edge of every cell. Each boundary edge lies in one
cell, and for $n_h\ge2$ a cell has at most two boundary edges (exactly two at the corners). Lemma~\ref{lem:trace} then
gives $\mathfrak b(\varphi,\varphi)\le C_h^2\,\mathfrak a_h(\varphi,\varphi)$.

(iii) Let $\mathcal Z_k$ be the span of the eigenfunctions of $\sigma_0,\dots,\sigma_k$, and take
$u\in\mathcal Z_k$ with $\mathfrak b(u,u)=1$. Put $\upsilon=\mathfrak a_h(u-\Pi_hu,u-\Pi_hu)$. By (i),
$\mathfrak a_h(\Pi_hu,\Pi_hu)=\int|\nabla u|^2-\upsilon\le\sigma_k-\upsilon$, and by (ii),
$\mathfrak b(\Pi_hu,\Pi_hu)^{1/2}\ge1-C_h\sqrt\upsilon$. If $C_h^2\sigma_k\ge1$ the claim is trivial, since
$x/(1+C_h^2x)<C_h^{-2}$. Otherwise $\Pi_h$ is injective on $\mathcal Z_k$, and
\[\frac{\mathfrak a_h(\Pi_hu,\Pi_hu)}{\mathfrak b(\Pi_hu,\Pi_hu)}\le\max_{0\le\upsilon\le\sigma_k}\frac{\sigma_k-\upsilon}{(1-C_h\sqrt\upsilon)^2}=\frac{\sigma_k}{1-C_h^2\sigma_k}.\]
The maximum is attained at $\sqrt\upsilon=C_h\sigma_k$. The min--max principle on the $(k+1)$-dimensional space
$\Pi_h\mathcal Z_k$ gives $\sigma_{h,k}\le\sigma_k/(1-C_h^2\sigma_k)$, which rearranges to the claim.
\end{proof}

\begin{proposition}\label{prop:cr}
$\sigma_4(E)>\clsigmaFourLo>2\sqrt3$.
\end{proposition}
\begin{proof}
Take $n_h=\clcrNh$, which gives $\clcrDofs$ degrees of freedom and $C_h^2\approx\clcrChSq$. Let $\mathbf S_h$ and
$\mathbf B_h$ be the matrices of $\mathfrak a_h$ and $\mathfrak b$ in the edge basis of $\mathcal C_h$. On an
equilateral cell the element stiffness matrix is
\[\frac{2}{\sqrt3}\begin{pmatrix}2&-1&-1\\-1&2&-1\\-1&-1&2\end{pmatrix}.\]
On a boundary edge, the basis function of that edge is $\equiv1$ and the other two are $\pm$ the linear function
rising from $-1$ to $1$ along the edge, so the element boundary-mass entries are $h$, $h/3$ and $-h/3$. Let $\mathbf O$ be a fixed matrix with binary64 entries (computed once as approximate
eigenvectors of $\mathbf S_h-5\mathbf B_h$). In ball arithmetic we check two things:
\begin{itemize}
\item the Gershgorin discs of $\mathbf O^T\mathbf O$ exclude $0$, so $\mathbf O$ is nonsingular;
\item the Gershgorin discs of $\mathbf O^T(\mathbf S_h-5\mathbf B_h)\mathbf O$ consist of $\clcrNeg$ discs in
$(-\infty,0)$ and $\clcrPos$ discs in $(0,\infty)$.
\end{itemize}
Gershgorin's counting theorem and Sylvester's law of inertia give exactly four negative eigenvalues of
$\mathbf S_h-5\mathbf B_h$. If $\sigma_{h,4}<5$, the pencil would be negative definite at this shift on
the span of five discrete eigenfunctions, including the constant. Thus $\sigma_{h,4}\ge5$, and
Lemma~\ref{lem:liu} gives $\sigma_4(E)\ge5/(1+5C_h^2)>\clsigmaFourLo$.
\end{proof}

\begin{proposition}\label{prop:eq}
For the side-one equilateral triangle,
\[
\sigma_1(E)=\sigma_2(E)=:\sigma^*\in J,\qquad \sigma_3(E)=2\sqrt3<\sigma_4(E).
\]
In particular the third positive eigenvalue is simple and
\[\Sigma(E)=3\sigma^*\in[\clSigmaElo,\,\clSigmaEhi].\]
\end{proposition}
\begin{proof}
Use Cartesian coordinates centred at the centroid and put $r=1/(2\sqrt3)$. The harmonic polynomial
$u=\Im((x+iy)^3)+2r^3$ satisfies $\partial_\nu u=-3x^2+3r^2=u/r$ on $y=-r$.
Rotation by $2\pi/3$ preserves $u$, so the same boundary condition holds on every edge.
Thus $1/r$ is an eigenvalue. The inclusion interval $J$ lies strictly between $0$ and $1/r$ and
contains at least two eigenvalues counted with multiplicity. Together with $0$ and $1/r$ these
account, with multiplicity, for all eigenvalues below the lower bound in Proposition~\ref{prop:cr}. Hence $J$ contains precisely
$\sigma_1,\sigma_2$, and $\sigma_3=1/r$ is simple. Take $\sigma^{\rm lo}=2\sqrt3$ in Lemma~\ref{lem:enclosure}.
Let $P$ be the spectral projection onto $J$. It commutes with the
$D_3$-action, so $(Pu_1,Pu_2)$ is again an equivariant pair. Since $\tilde\sigma>0$, Lemma~\ref{lem:enclosure} gives
$\|(I-P)u_a\|\le\eta\|u_a\|/(\sigma^{\rm lo}-\tilde\sigma)<10^{-12}\|u_a\|$; since $u_1\perp u_2$ with equal
norms, $Pu_1,Pu_2$ are linearly independent. They therefore span the two-dimensional range of $P$, which
carries the irreducible representation of $D_3$. By Schur's lemma $\Lam$ acts on this range as a scalar.
\end{proof}

\section{Local maximality}\label{sec:local}

Let $\mathcal U$ be the eigenspace of $\sigma^*$, with basis $u^*_a=Pu_a$ ($a=1,2$). For
$T=T(p,q)$, $p=\frac12+\rho\cos\psi$, $q=\frac{\sqrt3}{2}+\rho\sin\psi$, we have $T=A(E)$ with
\[A=I+\rho\begin{pmatrix}0&\alpha\\0&\beta\end{pmatrix},\quad \alpha=\tfrac{2\cos\psi}{\sqrt3},\ \beta=\tfrac{2\sin\psi}{\sqrt3},\quad d:=\det A=1+\rho\beta .\]
Number the edges $e_0=[(0,0),(1,0)]$, $e_1=[(1,0),\text{apex}]$ and $e_2=[\text{apex},(0,0)]$. Then
$s_0=1$, while
\[s_1=\bigl|\bigl(\rho\cos\psi-\tfrac12,\ \tfrac{\sqrt3}2+\rho\sin\psi\bigr)\bigr|,\qquad
s_2=\bigl|\bigl(\rho\cos\psi+\tfrac12,\ \tfrac{\sqrt3}2+\rho\sin\psi\bigr)\bigr|\]
are the side lengths of $T$, and $L_T=1+s_1+s_2$.

The \emph{moments} of $\mathcal U$ are the following $2\times2$ matrices and vectors, indexed by $a,b$: the
energy matrices $Q_{ij}=(\int_E\partial_iu^*_a\partial_ju^*_b)$, the edge Gram matrices
$g^e=(\int_eu^*_au^*_b)$, and the edge-mean vectors $\mu^e=(\int_eu^*_a)$. They satisfy two exact identities:
\begin{equation}\label{eq:exact}
\textstyle\sum_e\mu^e=0,\qquad Q_{11}+Q_{22}=\sigma^*\sum_eg^e .
\end{equation}
The first holds because eigenfunctions of $\Lam$ with $\sigma^*\neq0$ are orthogonal to constants. The
second is the weak eigen-equation on $\mathcal U$.

\begin{lemma}[Variational comparison]\label{lem:crit}
Let a trial space on a polygon $T$ have energy matrix $N$ and positive definite boundary-variance
matrix $D$. For a reference value $S$ put $X=(S/L(T))D-N$. Then
\[
S-\Sigma(T)\ge L(T)\max_{\xi\ne0}\frac{\xi^TX\xi}{\xi^TD\xi}.
\]
In particular, if $X-tD$ has a positive direction, the deficit exceeds $L(T)t$.
\end{lemma}
\begin{proof}
Apply~\eqref{eq:var} to each trial function and subtract the resulting bound from $S$.
\end{proof}

Apply the lemma with $S=3\sigma^*$, $N=\mathcal N_A$ and $D=\mathcal D_A$ on $\mathcal U$.
At $\rho=0$ the eigen-equation gives $X=0$. The following factorisation isolates the leading term.
\begin{lemma}[$X=\rho Y$]\label{lem:XrhoY}
Let $\lambda_e=(s_e-1)/\rho$. Then $\lambda_e=(\kappa_e+\rho)/(1+s_e)$, where
$\kappa_1=-\cos\psi+\sqrt3\sin\psi$ and $\kappa_2=\cos\psi+\sqrt3\sin\psi$. Assume~\eqref{eq:exact}. Then
$X=\rho Y$, where
\[Y=\sigma^*\Bigl(\sum_e\omega_eg^e-\frac{3\rho}{L_T^2}vv^T\Bigr)-\sum_{i,j}H_{ij}Q_{ij},\qquad
H=\begin{pmatrix}\beta+\rho\alpha^2/d&-\alpha/d\\-\alpha/d&-\beta/d\end{pmatrix},\]
with $\omega_0=-(\lambda_1+\lambda_2)/L_T$, $\omega_1=(2\lambda_1-\lambda_2)/L_T$,
$\omega_2=(2\lambda_2-\lambda_1)/L_T$ and $v=\lambda_1\mu^1+\lambda_2\mu^2$.
\end{lemma}
\begin{proof}
First, $s_e^2=1+\rho\kappa_e+\rho^2$, so $s_e-1=\rho(\kappa_e+\rho)/(1+s_e)$.

Next, $A^{-1}=\bigl(\begin{smallmatrix}1&-\rho\alpha/d\\0&1/d\end{smallmatrix}\bigr)$, which gives
$dG-I=\rho H$.

For the boundary terms, $\sum_es_e\mu^e=\sum_e(s_e-1)\mu^e=\rho v$ by~\eqref{eq:exact}. Also
$(3/L_T)s_e-1=(3s_e-L_T)/L_T=\rho\,\omega_e$, using $s_0=1$.

Hence $\sigma^*(3/L_T)\mathcal D_A-\sigma^*\sum_eg^e=\rho\,\sigma^*(\sum\omega_eg^e-3\rho vv^T/L_T^2)$, and
$\mathcal N_A-(Q_{11}+Q_{22})=\rho\sum H_{ij}Q_{ij}$. Subtracting these two equations and
using~\eqref{eq:exact} gives $X=\rho Y$.
\end{proof}

\begin{lemma}[Moment enclosures]\label{lem:moments}
Let $\delta_a=u_a-u^*_a$ and $\gamma=2\sqrt3-\tilde\sigma$. Then
\[\|\delta_a\|_{L^2(\partial E)}\le\eps_\partial:=\eta\|u_a\|_{\partial E}/\gamma,\qquad
\|\nabla\delta_a\|_{L^2(E)}\le\eps_\nabla:=\sqrt{2\sqrt3}\;\eta\|u_a\|_{\partial E}/\gamma .\]
Consequently each moment of $\mathcal U$ differs from the corresponding moment of $(u_1,u_2)$ by at most
the following amounts:
\begin{itemize}
\item $\eps_\nabla(\|\nabla u_a\|+\|\nabla u_b\|+\eps_\nabla)$ for the entry $(a,b)$ of $Q_{ij}$;
\item $\eps_\partial(\|u_a\|_{\partial E}+\|u_b\|_{\partial E}+\eps_\partial)$ for the entry $(a,b)$ of $g^e$;
\item $|e|^{1/2}\eps_\partial$ for the entry $a$ of $\mu^e$.
\end{itemize}
\end{lemma}
\begin{proof}
Apply Lemma~\ref{lem:enclosure} to $f=u_a$, with $\|r\|\le\eta\|u_a\|$; its hypotheses were verified in the
proof of Proposition~\ref{prop:eq}. The moment bounds follow from Cauchy--Schwarz.
\end{proof}

Here $\|u_a\|^2_{\partial E}\approx\clgramE$, so $\eps_\partial\le\clepsBdryMax$ and
$\eps_\nabla\le\clepsGradMax$. We compute the moments of $u_a$ exactly by Green's formula
$\int_E\mathfrak g\dd A=\frac1{2i}\oint_{\partial E}\bar z\,\mathfrak g\dd z$ and
$\int_E\mathfrak f_a'\overline{\mathfrak f_b'}\dd A=\frac1{2i}\oint_{\partial E}\overline{\mathfrak f_b}\mathfrak f_a'\dd z$ for holomorphic $\mathfrak g$. This turns every area integral into exact
polynomial line integrals.

\begin{proposition}\label{prop:local}
Let $\rho_1$ be the binary64 number nearest $0.6$. For $0<\rho\le\rho_1$ and every direction,
$\Sigma(E)-\Sigma(T(p,q))>\rho/2$.
\end{proposition}
\begin{proof}
We cover $[0,\rho_1]\times[0,\psi_{\max}]$ by $\cllocalBoxes$ closed boxes, where
$\psi_{\max}$ is the binary64 number nearest $6.3$ and exceeds $2\pi$. The boxes lie in the rectangle,
have disjoint interiors, and their exact rational areas sum to its area. They cover the rectangle since
a nonempty relatively open complement would have positive area. We evaluate $Y$ and $\mathcal D_A$ on
every box using Lemma~\ref{lem:moments}. For $Z=2Y-\mathcal D_A/2$ the checks give
\[
(\mathcal D_A)_{11}>0,\qquad \det\mathcal D_A>0,\qquad \tr Z>0\ \text{ or }\ \det Z<0 .
\]
Thus $Y-\mathcal D_A/4$ has a positive direction. Since $X=\rho Y$, Lemma~\ref{lem:crit}
gives a deficit greater than $L_T\rho/4$. The other two sides have total length greater than the unit
base, so $L_T>2$.
\end{proof}

\begin{theorem}[Splitting of the first pair]\label{thm:cone}
For a fixed unit base let the apex move from $(1/2,\sqrt3/2)$ by
$h=\rho(\cos\psi,\sin\psi)$. As $\rho\downarrow0$, uniformly in $\psi$,
\[
 L_T\sigma_1(T)=\Sigma(E)-c\rho+O(\rho^2),\qquad
 L_T\sigma_2(T)=\Sigma(E)+c\rho+O(\rho^2),
\]
where the exact constant satisfies $\clconeLo<c<\clconeHi$.
\end{theorem}
\begin{proof}
Work on $H=H^1(E)/\R$ with the Dirichlet inner product. On this fixed space let
$a_h=L_T\mathcal N_A$ and $b_h=\mathcal D_A$ be the pulled-back forms of~\eqref{eq:transplant},
now acting on all of $H$. Their representing operators $P_h,B_h$ are real analytic in the Cartesian
components of $h$: the energy coefficients are rational in $A$, and the edge lengths and perimeter
are positive near $h=0$. The operator $P_h$ is uniformly positive there. The boundary variance is
unchanged by adding constants and the trace is compact into $L^2(\partial E)$, so $B_h$ is compact.
Positive eigenvalues of $K_h=P_h^{-1/2}B_hP_h^{-1/2}$ are the reciprocals of the normalised positive
Steklov eigenvalues, by the weak equation on the quotient (choose the representative of zero
weighted boundary mean). The inverse square root is analytic near $P_0=3I$ by its power series.

At $h=0$ the top eigenvalue of $K_0$ is isolated and double by Proposition~\ref{prop:eq}.
The resolvent on a small enclosing circle has a convergent Neumann series for small $|h|$.
The Riesz projection defined by this contour is real analytic and has rank two by continuity. Project a basis at $0$
and apply $P_h^{-1/2}$ to obtain an analytic frame of the corresponding pencil subspace.
Multiplication by the inverse square root of its positive $b_h$-Gram matrix makes the frame
$b_h$-orthonormal. In this frame the matrix $M(h)$ of $a_h$ is symmetric and real analytic, with
eigenvalues $L_T\sigma_1(T), L_T\sigma_2(T)$. Write
\[
M(h)=\Sigma(E)I+F(h)+O(|h|^2),
\]
where $F$ is linear and the remainder is in matrix norm, uniformly in direction.

For a unit vector $e=(\cos\psi,\sin\psi)$ the matrix $F(e)$ represents
$\partial_e(a_h-\Sigma(E)b_h)|_0$ on the unperturbed eigenspace in a $b_0$-orthonormal basis.
Indeed, differentiating the moving frame introduces two further terms, each zero by the weak
eigen-equation against every vector of $H$. In the original basis $u_a^*$ the same derivative is
$-3Y(0,\psi)$, since $a_h-\Sigma(E)b_h=-L_T\rho Y$ on that fixed space. Consequently
\[
 \det F(e)=9\,\frac{\det Y(0,\psi)}{\det\mathcal D_I}.
\]

To see why the slope is independent of $e$, identify the apex with $z=p+iq$.
Cyclic relabelling of the vertices followed by a similarity replaces $z$ by $f(z)=1/(1-z)$.
It fixes $z_*=e^{i\pi/3}$ and $f'(z_*)=e^{2\pi i/3}$. Thus the two normalised eigenvalues agree as
an unordered pair at $z$ and $f(z)$. Both the trace and the determinant of $M(h)-\Sigma(E)I$
are therefore invariant under $h\mapsto f(z_*+h)-z_*$. Comparing the terms of degree one in the
trace and degree two in the determinant shows that $\tr F(h)$ and $\det F(h)$ are invariant under
rotation by $2\pi/3$. The only invariant linear form is zero, and every invariant real quadratic form
is a multiple of $|h|^2$ (its symmetric matrix commutes with that rotation).
Hence $F(h)$ has eigenvalues $\pm c|h|$, where
\[
 c=3\sqrt{-\det Y(0,0)/\det\mathcal D_I}.
\]
Ball evaluation using the enclosed moments of the projected eigenfunctions gives the stated positive
enclosure for $c$. The min--max principle bounds the change in each ordered eigenvalue by the norm
of the matrix remainder, proving both expansions.
\end{proof}

\section{The far region and the proof of the theorem}\label{sec:far}

For $0\le p\le\frac12$ put
\[q_{\max}(p)=\min\Bigl(\sqrt{2p-p^2},\ \tfrac{\sqrt3}{2}-\sqrt{\rho_1^2-(p-\tfrac12)^2}\Bigr).\]

\begin{lemma}\label{lem:cos}
\textup{(i)} $F\cap\{\rho\ge\rho_1\}=\{(p,q):0<p\le\frac12,\ 0<q\le q_{\max}(p)\}$.

\textup{(ii)} For $(p,q)\in F$, with $\ell_1=|(p,q)|$ and $\ell_2=|(1-p,q)|$,
\[\Sigma(T(p,q))\le\pi^2q\,\Xi(p)\,(1+\ell_1+\ell_2),\qquad \Xi(p)=\mathcal I(p)+\mathcal I(1-p),\]
where $\mathcal I(t)=\frac1t\int_0^tx\sin^2(\pi x)\dd x=\frac t4-\frac{\sin2\pi t}{4\pi}+\frac{\sin(\pi t)\sinc(\pi t)}{4\pi}$
and $\sinc x=\sin x/x$.
\end{lemma}
\begin{proof}
(i) For $(p,q)\in F$ we have $q\le\sqrt{2p-p^2}\le\frac{\sqrt3}2$, and $\rho\ge\rho_1$ means
$(q-\frac{\sqrt3}2)^2\ge\rho_1^2-(p-\frac12)^2$, where the right side is positive because
$|p-\frac12|\le\frac12<\rho_1$. Since $q<\frac{\sqrt3}2+\sqrt{\rho_1^2-(p-\frac12)^2}$, this is equivalent to
$q\le\frac{\sqrt3}2-\sqrt{\rho_1^2-(p-\frac12)^2}$.

(ii) Take $u=\cos(\pi x)$ in~\eqref{eq:var}. The vertical section of $T$ at abscissa $x$ has length $qx/p$
for $x\le p$ and $q(1-x)/(1-p)$ for $x\ge p$, so the substitution $x\mapsto1-x$ on $[p,1]$ gives
\[\int_T|\nabla u|^2=\pi^2q\Bigl(\frac1p\int_0^px\sin^2\pi x\dd x+\frac1{1-p}\int_0^{1-p}x\sin^2\pi x\dd x\Bigr)=\pi^2q\,\Xi(p).\]
The antiderivative $\frac{x^2}4-\frac{x\sin2\pi x}{4\pi}-\frac{\cos2\pi x}{8\pi^2}$ and the identity
$(1-\cos2\pi t)/t=2\pi\sin(\pi t)\sinc(\pi t)$ give the closed form of $\mathcal I$. The denominator of~\eqref{eq:var}
is $\min_\zeta\int_{\partial T}(u-\zeta)^2$. On the side from $(0,0)$ to $(p,q)$ arc length is
$(\ell_1/p)\dd x$ with $\ell_1/p\ge1$, and similarly on the third side. So for every constant $\zeta$,
\[\int_{\partial T}(u-\zeta)^2\ge2\int_0^1(\cos\pi x-\zeta)^2\dd x=1+2\zeta^2\ge1.\]
Hence $\Sigma(T)\le L\int_T|\nabla u|^2$ with $L=1+\ell_1+\ell_2$.
\end{proof}

\begin{proposition}\label{prop:far}
For every $(p,q)\in F$ with $\rho\ge\rho_1$, $\Sigma(T(p,q))\le\clfarBound<\Sigma(E)$.
\end{proposition}
\begin{proof}
By Lemma~\ref{lem:cos}(i) these are the points with $0<p\le\frac12$, $0<q\le q_{\max}(p)$. The bound of
Lemma~\ref{lem:cos}(ii) is nondecreasing in $q$, and $\Xi\ge0$. On an interval $p_-\le p\le p_+$, the
first entry of $q_{\max}$ is increasing in $p$ and the second is decreasing, so
\[q_{\max}(p)\le q_+:=\min\Bigl(\sqrt{2p_+-p_+^2},\ \tfrac{\sqrt3}2-\sqrt{\rho_1^2-(p_--\tfrac12)^2}\Bigr).\]
Moreover $\ell_1\le(p_+^2+q_+^2)^{1/2}$ and $\ell_2\le((1-p_-)^2+q_+^2)^{1/2}$. Hence on that interval
\[\Sigma(T)\le\pi^2q_+\,\sup_{[p_-,p_+]}\Xi\,\Bigl(1+\sqrt{p_+^2+q_+^2}+\sqrt{(1-p_-)^2+q_+^2}\Bigr).\]
We evaluate the right side in ball arithmetic on the $16$ intervals $[\frac k{32},\frac{k+1}{32}]$,
$0\le k\le15$, which cover $[0,\frac12]$. Each upper endpoint is below $\clfarBound$.
\end{proof}

\subsection*{Proof of Theorem~\ref{thm:main}}
Use the normalisation of Lemma~\ref{lem:param}. The deficit vanishes at $\rho=0$ and exceeds
$\rho/2$ for $0<\rho\le\rho_1$ by Proposition~\ref{prop:local}. For $\rho\ge\rho_1$,
Proposition~\ref{prop:far} gives a deficit at least $\clSigmaElo-\clfarBound>1/2$.
Throughout $F$ we have $0<p\le1/2$ and $0<q\le\sqrt3/2$, so $\rho<1$.
Finally $|(p,q)|=\ell_{\min}/\ell_{\max}$ in these coordinates, and the reverse triangle
inequality gives $1-|(p,q)|\le\rho$. This proves both deficit bounds and the equality case.
The extremal value is enclosed in Proposition~\ref{prop:eq}.
\qed

\begin{corollary}[Fixed-profile thin limit]\label{cor:thin}
Fix $0<p<1$, and set $\Phi_p(x)=\min\{x/p,(1-x)/(1-p)\}$ on $[0,1]$. Define
\[
 \theta(p)=\inf_{\substack{0\ne v\in H^1(0,1)\\ \int_0^1v=0}}
 \frac{\int_0^1\Phi_p(x)|v'(x)|^2\dd x}{2\int_0^1|v(x)|^2\dd x}.
\]
Then $\sigma_1(T(p,q))=q\theta(p)+o(q)$ and $\Sigma(T(p,q))=2q\theta(p)+o(q)$ as $q\downarrow0$.
\end{corollary}
\begin{proof}
In Henrot and Michetti's thin-domain lemma~\cite[Lemma~3.2]{HM} take the upper profile $\Phi_p$,
the lower profile zero, and thickness $q$. Their lower-growth hypothesis holds since
$\Phi_p(x)\ge x(1-x)$. Their first one-dimensional eigenvalue is $2\theta(p)$, and their
asymptotic formula multiplies it by $q/2$. Now use $L(T(p,q))\to2$.
\end{proof}

\section{Implementation}\label{sec:cert}

All certified steps use Arb~\cite{Arb} (python-flint~0.9.0), at $64$--$200$ bits.

\textbf{Certificate data.} Floating-point arithmetic is used only to produce candidate data, which are then
fixed and treated as exact binary numbers: the coefficients $c_k$ and the congruence matrix $\mathbf O$. These
are committed to the repository, and the certified steps read only them. The shift $\tilde\sigma$ is computed
from the $c_k$ in ball arithmetic and rounded to binary64. Hence the certified results do not depend on the
linear-algebra library or the processor; regenerating the data is a separate, optional step. Every number of
this paper that a computation checks is kept in a single claims file. The checking script compares each
computed enclosure with these claims in exact rational arithmetic, and the paper's macros are generated from
the same file.

\textbf{Boxes.} Box endpoints are exact binary numbers. A parameter box $[x_0,x_1]$ enters as the Arb union of its two exact endpoints, which contains
the whole interval. Square roots of nonnegative quantities whose balls straddle $0$ are clamped to
$[0,\sqrt{\cdot}]$. In the code, $\rho_1$ is the binary double nearest $0.6$, which is
$0.59999999999999997779\ldots$, slightly below $0.6$. The local and the far computations use this same value,
so the two regions meet exactly. The tiling of the local boxes is checked with exact float comparisons and an
exact rational area sum.

\textbf{Consistency checks} (not part of the proof).
\begin{itemize}
\item The Green-formula moments agree with brute-force quadrature on random harmonic polynomials to
$5\times10^{-14}$.
\item The transplant bound of \S\ref{sec:local}, and the full Rayleigh quotient of $\cos(\pi x)$, agree with
direct quadrature on $T$ to about $10^{-14}$.
\item At $20$ sample triangles both bounds lie above the floating-point Ritz value of $\Sigma$, as upper
bounds must.
\item The Crouzeix--Raviart discrete values are $\sigma_{h,1}=\clcrSigmaEight$ for $n_h=8$ and $\clcrSigmaSixteen$ for
$n_h=16$, increasing towards $\sigma^*=1.29082$.
\end{itemize}

\textbf{Code names.} Some identifiers differ from the notation here. In the Python code, $\tilde\sigma$ is
\texttt{sig}, the edge means $\mu^e$ are \texttt{mm}, $\lambda_1,\lambda_2$ are \texttt{d1}, \texttt{d2},
the $\omega_e$ are \texttt{ce}, and $\mathcal D_A$ is \texttt{D}. In Lean, the $\mu^e$ are \texttt{m0x},
\dots, \texttt{m2y}, the $\omega_e$ are \texttt{c0}, \texttt{c1}, \texttt{c2}, and $\mathcal N_A,\mathcal D_A$
are \texttt{NA}, \texttt{DA}.

\textbf{Lean.} A Lean~4/Mathlib development formalises the elementary steps:
\begin{itemize}
\item the triangle isoperimetric inequality in Heron form, with its equality case, and the deduction of
Corollary~\ref{cor:area};
\item the algebra of Lemma~\ref{lem:XrhoY};
\item the integral defining $\mathcal I$ in Lemma~\ref{lem:cos}, and the two edge integrals of $\cos(\pi x)$ used in
the consistency checks.
\end{itemize}
It uses no \texttt{sorry}, and its only axioms are \texttt{propext}, \texttt{Classical.choice} and
\texttt{Quot.sound}. It does not formalise the spectral theory, the trace and Payne--Weinberger
inequalities, or the ball-arithmetic computations.

\textbf{Availability.} The code, the logs, the box certificates and the Lean files are in the accompanying
repository\footnote{\url{https://github.com/jasondavies/steklov-triangles}.}. The single command \texttt{scripts/verify.py} reruns all of them and writes a PASS/FAIL
summary.

\section{Numerical observations (not rigorous)}\label{sec:numerics}

Except for the certified upper bounds in Table~\ref{tab:values}, the statements in this section are
numerical. They use the harmonic-polynomial Ritz method, with trial space $\{\Re w^k,\Im w^k\}_{k\le n}$ and
all integrals computed on the boundary.

\begin{observation}[Global picture]
On a $25\times30$ grid over $F$ (double precision, $n=30$) the largest value of $\Sigma$ is attained at the
equilateral apex. Table~\ref{tab:values} lists selected values, sorted by $\Sigma$, together with the regular
$N$-gons for $N=4,5,6$. They are computed as follows:
\begin{itemize}
\item the Ritz matrices are exact polynomial line integrals, evaluated in ball arithmetic;
\item the eigenproblem is solved in $50$- and $80$-digit arithmetic with the full basis of $2n$
directions, for $n=30,40,50$;
\item the estimate is the value for $n=50$;
\item its error is estimated (not rigorously) as twice the largest of the last change in $n$, a geometric
and a power-law extrapolation of the remaining changes, and the change between the two precisions; the
estimate is printed only to the digits this error justifies;
\item the certified upper bound is the Rayleigh quotient of the fixed $n=50$ trial function, evaluated in
ball arithmetic on the exact polygon.
\end{itemize}
Convergence is super-geometric for $E$, the square, and the right-angled triangles. It is slowest for the
polygons with an angle of $108^\circ$ or $120^\circ$ (the $30$--$30$--$120$ triangle, the pentagon and the
hexagon), whose corners produce the strongest singularities.

\begin{table}[htbp]
\centering
\begin{tabular}{lll}
\toprule polygon & certified upper bound for $\Sigma$ & estimate of $\Sigma$\\\midrule
$60$--$60$--$60$ & $\clTequiUp$ & $\clTequiEst\pm\clTequiErr$\\
$50$--$50$--$80$ & $\clTfiftyUp$ & $\clTfiftyEst\pm\clTfiftyErr$\\
$70$--$70$--$40$ & $\clTseventyUp$ & $\clTseventyEst\pm\clTseventyErr$\\
$40$--$60$--$80$ & $\clTfortyUp$ & $\clTfortyEst\pm\clTfortyErr$\\
$45$--$45$--$90$ & $\clTrightisoUp$ & $\clTrightisoEst\pm\clTrightisoErr$\\
$30$--$60$--$90$ & $\clThalfeqUp$ & $\clThalfeqEst\pm\clThalfeqErr$\\
$30$--$30$--$120$ & $\clTobtuseUp$ & $\clTobtuseEst\pm\clTobtuseErr$\\
\midrule
square & $\clTsquareUp$ & $\clTsquareEst\pm\clTsquareErr$\\
regular pentagon & $\clTpentagonUp$ & $\clTpentagonEst\pm\clTpentagonErr$\\
regular hexagon & $\clThexagonUp$ & $\clThexagonEst\pm\clThexagonErr$\\
\bottomrule
\end{tabular}
\caption{Values of $\Sigma=\sigma_1L$. The $\pm$ values are heuristic error estimates.
Rounding the central estimate to the displayed precision can place it slightly above the certified upper bound.}\label{tab:values}
\end{table}

\smallskip
The regular $N$-gon values increase with $N$, consistent with~\cite{CGHLY}.
\end{observation}

\begin{observation}[Thin triangles]
For the limiting endpoint profile $\Phi_0(x)=1-x$ in the quotient of Corollary~\ref{cor:thin}, the
Euler--Lagrange equation $((1-x)u')'+2\theta u=0$ gives $\theta(0)=j_{1,1}^2/8\approx\clthetaZero$. Numerically
$\theta(0.05)=\clthetaA$, $\theta(0.25)=\clthetaB$ and $\theta(\frac12)=\clthetaC$. The double-precision Ritz values ($n=40$)
of $\Sigma/(2q)$ at $q=0.01$ agree with these to $\clthinAgree$ relative. In particular
$\Sigma(T)\approx2.89\,\vartheta(T)$ for thin isosceles triangles, compared with the rigorous constant $45.6$
of~\cite{DGMRV}.
\end{observation}

\section{Remarks}

Lemma~\ref{lem:crit} applies to any polygonal trial space for which the energy and boundary-variance forms
can be enclosed. At a multiple reference eigenvalue the spectral gap controls the projection error,
and the identity $X=\rho Y$ makes verification of a linear deficit possible up to $\rho=0$.

For $n\ge4$ the affine group no longer acts transitively on similarity classes of $n$-gons, so one needs
a richer family of transplantations, for example piecewise affine or Schwarz--Christoffel maps. Parallelograms form a two-parameter affine family of the square. General quadrilaterals
have four parameters modulo similarities, so a first-order comparison of a double eigenvalue would also
have to address deformations in the kernel of the linear map into symmetric $2\times2$ matrices.

\section*{Declaration of AI-generated research and authorship}
This manuscript is the product of an AI-driven research process. During a
broader investigation of open problems from the AIM problem lists, initiated
by the human author, an AI system identified this problem as a candidate,
found the proof, and wrote the code for the certified computations. The author
then directed AI systems to package the result, formalise parts of it in Lean
and prepare this manuscript. The AI systems developed the results and their
presentation, including the mathematical exposition, the computations, the
formalisation, and the organisation and writing of the text; the human author
wrote no part of the manuscript. The system was Claude Opus~5.5 (Anthropic),
used through Claude Code. OpenAI Codex subsequently revised the proof and computational certificates for
stability, spectral isolation and splitting, and incorporated the thin-domain corollary.

The proof is computer-assisted and is reproduced by \texttt{scripts/verify.py}; the Lean development covers
only the elementary steps listed in \S\ref{sec:cert}.

The human author has not read this manuscript or checked its arguments. At
the time of writing, the mathematics, the code and the formalisation have
been checked only by AI systems and by the computer verification described
above; no independent human expert has reviewed them. The author takes
responsibility for the accuracy of this disclosure.

\raggedright
\begin{thebibliography}{99}
\bibitem{AimPL} AIM Problem Lists, \emph{Steklov eigenproblems}, Section~4 (Shape optimisation for Steklov), Problem~4.6 (\'E.~Oudet), \url{http://aimpl.org/stekloveigen/4}.
\bibitem{Arb} F.~Johansson, \emph{Arb: efficient arbitrary-precision midpoint-radius interval arithmetic}, IEEE Trans.\ Comput.\ \textbf{66} (2017), 1281--1292.
\href{https://doi.org/10.1109/TC.2017.2690633}{doi:10.1109/TC.2017.2690633}.
\bibitem{Beb} M.~Bebendorf, \emph{A note on the Poincar\'e inequality for convex domains}, Z.~Anal.\ Anwend.\ \textbf{22} (2003), 751--756.
\href{https://doi.org/10.4171/ZAA/1170}{doi:10.4171/ZAA/1170}.
\bibitem{CGHLY} Z.~Cheng, C.~Gui, Y.~Hu, Q.~Li, R.~Yao, \emph{Global strict monotonicity and asymptotics of the first Steklov eigenvalue for regular polygons}, arXiv:2603.25116v2 (2026).
\href{https://doi.org/10.48550/arXiv.2603.25116}{doi:10.48550/arXiv.2603.25116}.
\bibitem{CR} M.~Crouzeix, P.-A.~Raviart, \emph{Conforming and nonconforming finite element methods for solving the stationary Stokes equations~I}, RAIRO Anal.\ Num\'er.\ \textbf{7} (1973), 33--75.
\href{https://doi.org/10.1051/m2an/197307R300331}{doi:10.1051/m2an/197307R300331}.
\bibitem{DK} C.~Davis, W.~M.~Kahan, \emph{The rotation of eigenvectors by a perturbation.~III}, SIAM J.\ Numer.\ Anal.\ \textbf{7} (1970), 1--46.
\href{https://doi.org/10.1137/0707001}{doi:10.1137/0707001}.
\bibitem{DNS} S.~Dominguez, N.~Nigam, B.~Shahriari, \emph{A combined finite element and Bayesian optimization framework for shape optimization in spectral geometry}, Comput.\ Math.\ Appl.\ \textbf{74} (2017), no.~11, 2874--2896.
\href{https://doi.org/10.1016/j.camwa.2017.08.044}{doi:10.1016/j.camwa.2017.08.044}.
\bibitem{DGMRV} E.~B.~Dryden, C.~Gordon, J.~Moreno, J.~Rowlett, C.~Villegas-Blas, \emph{The Steklov spectrum of convex polygonal domains~I: spectral finiteness}, J.\ Geom.\ Anal.\ \textbf{35} (2025), no.~3, Paper No.~91.
\href{https://doi.org/10.1007/s12220-025-01922-8}{doi:10.1007/s12220-025-01922-8}.
\bibitem{GLPS} A.~Girouard, J.~Lagac\'e, I.~Polterovich, A.~Savo, \emph{The Steklov spectrum of cuboids}, Mathematika \textbf{65} (2019), no.~2, 272--310.
\href{https://doi.org/10.1112/S0025579318000414}{doi:10.1112/S0025579318000414}.
\bibitem{GP} A.~Girouard, I.~Polterovich, \emph{Spectral geometry of the Steklov problem}, J.\ Spectr.\ Theory \textbf{7} (2017), 321--359.
\href{https://doi.org/10.4171/JST/164}{doi:10.4171/JST/164}.
\bibitem{LS} R.~S.~Laugesen, B.~A.~Siudeja, \emph{Maximizing Neumann fundamental tones of triangles}, J.\ Math.\ Phys.\ \textbf{50} (2009), 112903.
\href{https://doi.org/10.1063/1.3246834}{doi:10.1063/1.3246834}.
\bibitem{HM} A.~Henrot and M.~Michetti, \emph{A comparison between Neumann and Steklov eigenvalues},
J. Spectr. Theory \textbf{12} (2022), 1405--1442, \href{https://doi.org/10.4171/JST/429}{doi:10.4171/JST/429}.
\bibitem{Nigam} N.~Nigam, Open Problem~8 in \emph{Open problems from the miniconference on sharp eigenvalue estimates
for partial differential operators} (2020), \href{https://publish.illinois.edu/eigenvalues2020/files/2020/04/Open-problems.pdf}{problem list}.
\bibitem{Liu} X.~Liu, \emph{A framework of verified eigenvalue bounds for self-adjoint differential operators}, Appl.\ Math.\ Comput.\ \textbf{267} (2015), 341--355.
\href{https://doi.org/10.1016/j.amc.2015.03.048}{doi:10.1016/j.amc.2015.03.048}.
\bibitem{PW} L.~E.~Payne, H.~F.~Weinberger, \emph{An optimal Poincar\'e inequality for convex domains}, Arch.\ Rational Mech.\ Anal.\ \textbf{5} (1960), 286--292.
\href{https://doi.org/10.1007/BF00252910}{doi:10.1007/BF00252910}.
\bibitem{Weinstock} R.~Weinstock, \emph{Inequalities for a classical eigenvalue problem}, J.\ Rational Mech.\ Anal.\ \textbf{3} (1954), 745--753.
\href{https://doi.org/10.1512/iumj.1954.3.53036}{doi:10.1512/iumj.1954.3.53036}.
\bibitem{YXL} C.~You, H.~Xie, X.~Liu, \emph{Guaranteed eigenvalue bounds for the Steklov eigenvalue problem}, SIAM J.\ Numer.\ Anal.\ \textbf{57} (2019), 1395--1410.
\href{https://doi.org/10.1137/18M1189592}{doi:10.1137/18M1189592}.
\end{thebibliography}

\end{document}
